\DocumentMetadata{
  pdfversion=2.0,
  pdfstandard=ua-2,
  lang=en-US,
  tagging=on,
  tagging-setup={math/setup=mathml-SE,tabsorder=structure}
}
\documentclass[11pt,letterpaper]{article}
\usepackage[margin=0.68in,includefoot]{geometry}
\usepackage{newcomputermodern}
\usepackage{amsmath,amscd,mathtools}
\usepackage{tikz,adjustbox}
\providecommand{\square}{\mdlgwhtsquare}
\usepackage[hidelinks]{hyperref}
\hypersetup{
  pdftitle={Topology First-Year Exam, January 2024, Part 1},
  pdfauthor={University of Florida Department of Mathematics},
  pdfsubject={Graduate mathematics examination},
  pdfdisplaydoctitle=true,
  bookmarksopen=true
}
\setcounter{secnumdepth}{0}
\setlength{\parindent}{0pt}
\setlength{\parskip}{0.28em}
\setlength{\emergencystretch}{3em}
\setlength{\leftmargini}{2.15em}
\setlength{\leftmarginii}{2.65em}
\setlength{\leftmarginiii}{3em}
\renewcommand{\labelenumi}{\textbf{\theenumi.}}
\pagestyle{empty}
\raggedbottom
\allowdisplaybreaks
\newenvironment{examproblems}[1]{%
  \begin{enumerate}%
  \setcounter{enumi}{#1}%
  \setlength{\topsep}{0.35em}%
  \setlength{\partopsep}{0pt}%
  \setlength{\itemsep}{0.48em}%
  \setlength{\parsep}{0.18em}%
}{\end{enumerate}}
\newenvironment{examsubparts}{%
  \begin{enumerate}%
  \setlength{\topsep}{0.18em}%
  \setlength{\partopsep}{0pt}%
  \setlength{\itemsep}{0.16em}%
  \setlength{\parsep}{0.1em}%
}{\end{enumerate}}
\newenvironment{examinstructions}{%
  \par\begingroup\itshape
}{%
  \par\endgroup\vspace{0.18em}
}
\makeatletter
\renewcommand\section{\@startsection{section}{1}{0pt}%
  {-1.25ex plus -.25ex minus -.1ex}%
  {0.55ex plus .1ex}%
  {\normalfont\large\bfseries}}
\renewcommand{\@maketitle}{%
  \newpage\null\vskip -1em%
  {\footnotesize\bfseries
    \href{https://www.ufl.edu/}{University of Florida}, \href{https://math.ufl.edu/}{Department of Mathematics}\par}%
  \vskip -0.5em%
  \tagstructbegin{tag=Title}%
    {\LARGE\bfseries\href{https://gma.math.ufl.edu/exams/topology-first-year/}{Topology First-Year Exam}, January 2024, Part 1\par}%
  \tagstructend%
  \vskip 0.25em%
  \tagmcbegin{artifact}%
    \hrule height 1pt%
  \tagmcend%
  \vskip 1em%
}
\makeatother
\title{Topology First-Year Exam, January 2024, Part 1}
\author{}
\date{}
\begin{document}
\maketitle
\thispagestyle{empty}
\begin{examinstructions}
For the first five problems, show all your work and support all statements.
\end{examinstructions}
\begin{examproblems}{0}
\item
Prove that \(\{0,1\}^{\omega}\) is not countable.
\item
Show that \(\mathbb{R}^{\omega}\) with the box topology is not connected.
\item
Assume that \(f:X\to Y\) is continuous and surjective.
\begin{examsubparts}
\item[\textnormal{(a)}]
If~\(X\) is Lindelöf (every open covering has a countable subcovering), show that~\(Y\) is also.
\item[\textnormal{(b)}]
If~\(X\) is separable (has a countable dense subset), show that~\(Y\) is also.
\end{examsubparts}
\item
Let \((X,d)\) be a compact metric space. Let \(f:X\to X\) and suppose there is some \(c<1\) with \(d(f(x),f(y))\leq c\cdot d(x,y)\) for all \(x,y\in X\). Prove there exists a unique \(x\in X\) with \(f(x)=x\).
\item
Let \((Y,d)\) be a complete metric space. Let~\(J\) be a set. Suppose \(d(y,y')\leq 1\) for all \(y,y'\in Y\), meaning a simpler definition for the uniform metric \(\bar{\rho}\) on \(Y^J\) is 
\[
\bar{\rho}(f,g)=\sup_{\alpha\in J}\{d(f(\alpha),g(\alpha))\}
\]
 for \(f,g:J\to Y\). Prove that \((Y^J,\bar{\rho})\) is a complete metric space.
\end{examproblems}
\begin{examinstructions}
Answer the following problems with complete definitions, complete statements, an example, or a short proof.
\end{examinstructions}
\begin{examproblems}{5}
\item
Define what it means for a topological space~\(X\) to be regular.
\item
Show that a closed subspace~\(C\) of a compact space~\(X\) is compact.
\item
A function \(f:\mathbb{Z}_+\to\{0,1\}\) is eventually zero if there is a positive integer~\(N\) such that \(f(n)=0\) for all \(n\geq N\). Prove that the set of all functions \(f:\mathbb{Z}_+\to\{0,1\}\) that are eventually zero is countable.
\item
Show that \([0,\infty)\) and~\(\mathbb{R}\) are not homeomorphic.
\item
Let \(x=(x_1,x_2,\ldots)\in\mathbb{R}^{\omega}\). Let \(\{x^{(n)}\}_{n\in\mathbb{Z}_+}\) be a sequence of points in \(\mathbb{R}^{\omega}\) such that for each coordinate \(i\in\mathbb{Z}_+\), the sequence \(\{x_i^{(n)}\}_{n\in\mathbb{Z}_+}\) converges to \(x_i\) in~\(\mathbb{R}\). Show that \(\{x^{(n)}\}_{n\in\mathbb{Z}_+}\) converges to~\(x\) if \(\mathbb{R}^{\omega}\) has the product topology.
\item
Give an example of a topological space that is connected but not path-connected.
\item
Define what it means for a topological space~\(X\) to be a Baire space.
\item
Show that if a topological space~\(X\) is Hausdorff, then a sequence in~\(X\) can converge to at most one point in~\(X\).
\item
Show that the function \(f:[0,2\pi)\to\{(x,y)\in\mathbb{R}^2\mid x^2+y^2=1\}\) from the half-open interval to the circle defined by \(f(t)=(\cos(t),\sin(t))\) is not a homeomorphism.
\item
Let~\(X\) be a locally compact Hausdorff space. Define the topology on \(Y=X\cup\{\infty\}\) that makes~\(Y\) the one-point compactification of~\(X\).
\end{examproblems}
\end{document}
